\documentclass[11pt]{article}
\usepackage[utf8]{inputenc}
\usepackage{amsmath,amsthm,amssymb,mathtools}
\usepackage{geometry}
\geometry{margin=2.5cm}
\usepackage[hidelinks]{hyperref}
\usepackage{booktabs}
\usepackage{tikz}

\newtheorem{theorem}{Theorem}[section]
\newtheorem{lemma}[theorem]{Lemma}
\newtheorem{corollary}[theorem]{Corollary}
\newtheorem{proposition}[theorem]{Proposition}

\theoremstyle{definition}
\newtheorem{definition}[theorem]{Definition}
\newtheorem{example}[theorem]{Example}

\theoremstyle{remark}
\newtheorem{remark}[theorem]{Remark}
\newtheorem{observation}[theorem]{Observation}

\newcommand{\floor}[1]{\lfloor #1 \rfloor}
\newcommand{\ceil}[1]{\lceil #1 \rceil}
\newcommand{\B}{\mathrm{B}}
\newcommand{\N}{\mathbb{N}}
\newcommand{\Z}{\mathbb{Z}}
\newcommand{\Q}{\mathbb{Q}}
\newcommand{\R}{\mathbb{R}}
\newcommand{\pFq}[2]{{}_{#1}F_{#2}}

\title{Closed-Form State Vectors for Lattice Paths\\
Below Irrational Staircases}
\author{Jan Popelka\thanks{Email: popojan@protonmail.com.
Code: \url{https://github.com/popojan/orbit}}}
\date{}

\begin{document}

\maketitle

\begin{abstract}
We derive explicit binomial formulas for the transfer matrix of
lattice paths constrained by the Beatty staircase
$\floor{x/\alpha}$, where $\alpha$ is irrational with continued
fraction $[w; a_1, a_2, \ldots]$.
The transfer matrix between consecutive convergents decomposes as
a Toeplitz matrix minus a correction given by a finite sum of
products of ballot numbers and binomial coefficients.
For the step width $w = 1$, this correction collapses to a single
binomial via the classical Rothe--Hagen identity.
For $w \geq 2$, the correction is an irreducible terminating
Saalsch\"utzian generalized hypergeometric function
$\pFq{2w+2}{2w+1}(1)$.
The formula is self-similar: it applies at every level of the
continued fraction hierarchy with only the local parameters
changing, and successive levels compose via a terminating
inclusion--exclusion expansion.
For $w = 1$ and quadratic irrationals, the level-to-level scaling
is governed by powers of the fundamental unit of
$\mathbb{Q}(\sqrt{D})$.
\end{abstract}

%======================================================================
\section{Introduction}\label{sec:intro}
%======================================================================

\subsection{Context}

Counting monotonic lattice paths subject to boundary constraints is
a classical problem in enumerative combinatorics, with origins in
Bertrand's ballot problem~\cite{Andre1887} and connections to
random walks, queuing theory, and algebraic
combinatorics~\cite{Krattenthaler2015}.

When the boundary is the line $y = qx/p$ for coprime $p, q$, the
number of paths from $(1,0)$ to $(p, j)$ staying weakly below the
boundary is given by a ballot-type formula involving a single
binomial coefficient~\cite{IrvingRattan2009}.
When the boundary is the Beatty staircase
$S(x) = \floor{x/\alpha}$ for an irrational $\alpha$, the path
count at convergent positions equals the ballot number
$\B(p,q) = \binom{p+q-1}{q}/p$, as shown
in~\cite{PopelkaBeattyBallot2026}.

But the ballot number gives only the \emph{total} path count at the
endpoint $(p, q)$.
The full \emph{state vector} --- the number of paths reaching each
intermediate height $j = 0, 1, \ldots, q$ at position~$p$ ---
contains richer information and is essential for composing blocks
across multiple continued fraction levels.

\subsection{Main results}

Let $\alpha = [w; a_1, a_2, \ldots]$ be irrational with convergents
$p_k/q_k$.
We study the \emph{block transfer matrix} $M$ that maps the state
vector at one convergent position to the next.
Our main contributions are:

\begin{enumerate}
\item \textbf{Uniform formula (Theorem~\ref{thm:uniform}).}
  For heights $j \leq q_1$, the state vector entry $v_j(p)$ at any
  semi-convergent position $p$ equals the rational ballot formula
  $v_j(p) = \frac{p - wj}{p}\binom{p+j-1}{j}$, with no correction
  needed.

\item \textbf{Block transfer decomposition
  (Theorem~\ref{thm:transfer}).}
  The full transfer matrix decomposes as $M = T - \Delta$, where
  $T$ is the Toeplitz matrix
  $T[j,s] = \binom{p_1 - 1 + j - s}{j - s}$ and $\Delta$ is
  a correction supported on rows $j \geq d_0$, given by
  %
  \begin{equation}\label{eq:R7-intro}
    \Delta[d_0{+}d,\, s] = \sum_{m=1}^{d+1}
      v_{d-m+1}(p_1 - wm)\;\binom{d_0 - 2 + m(w{+}1) - s}{mw - 1},
  \end{equation}
  %
  where $d_0$ is the input dimension and $d = 0, \ldots, q_1 - 1$.
  The decomposition is proven for all $q_1$ by a linearity and
  self-similarity argument: each rise injects a free-path count
  into the newly accessible row, which then propagates as a
  uniform-staircase path count in the residual block.

\item \textbf{Hypergeometric classification
  (Theorem~\ref{thm:hypergeom}).}
  The correction~\eqref{eq:R7-intro} is a terminating $1$-balanced
  (Saalsch\"utzian) generalized hypergeometric function
  $\pFq{2w+2}{2w+1}(1)$.
  For $w = 1$, it collapses via the Pfaff--Saalsch\"utz theorem to
  the single binomial $\binom{d_0 + p_1 + d - 1 - s}{d}$ (the
  Rothe--Hagen identity with parameter $b = 2$).
  For $w \geq 2$, no such collapse exists.

\item \textbf{Self-similar hierarchy (Theorem~\ref{thm:selfsim}).}
  Formula~\eqref{eq:R7-intro} applies at every continued fraction
  level with only $p_k$ and $d_{0,k}$ changing.
  Successive levels compose via a product of $a_{k+1}$ block
  matrices, whose inclusion--exclusion expansion terminates exactly
  at order $a_{k+1}$.

\item \textbf{Rotation decomposition (Theorem~\ref{thm:rotation}).}
  For $w = 1$, the block transfer at non-canonical Sturmian
  positions decomposes as
  $M_r = \operatorname{sturmianBlock}(d_0{+}r,\, p{-}r,\, q{-}r)
  \cdot \operatorname{riseProduct}(d_0,\, r)$,
  eliminating the level-$0$ ambiguity.
\end{enumerate}

\subsection{Relation to prior work}

Irving and Rattan~\cite{IrvingRattan2009} gave the uniform ballot
formula (our Result~1) for paths below a rational line $y = x/a$.
Banderier and Wallner~\cite{BanderierWallner2019} studied the
generating function approach, showing that individual heights under
rational staircases have complicated expressions but block sums
simplify (``ugly $+$ ugly $=$ nice'').
Firoozi, Jedwab, and Rattan~\cite{FirooziJedwabRattan2026} studied
flaw-counting under rational slopes using alternating sums
reminiscent of our correction structure.

Our contribution is the explicit closed-form correction
formula~\eqref{eq:R7-intro} for individual heights beyond the
uniform zone, its identification as a Saalsch\"utzian hypergeometric
function, and the self-similar hierarchical structure linking all
continued fraction levels.

\subsection{Outline}

Section~\ref{sec:prelim} fixes notation.
Section~\ref{sec:uniform} establishes the uniform formula.
Section~\ref{sec:transfer} proves the main decomposition theorem.
Section~\ref{sec:hypergeom} classifies the correction
hypergeometrically.
Section~\ref{sec:selfsim} develops the hierarchical structure.
Section~\ref{sec:rotation} treats the rotation decomposition.
Section~\ref{sec:wreduction} presents the $w$-reduction kernel.
Section~\ref{sec:scaling} discusses level scaling for quadratic
irrationals.
Section~\ref{sec:open} lists open problems.

%======================================================================
\section{Preliminaries}\label{sec:prelim}
%======================================================================

\subsection{Continued fractions}

Let $\alpha > 1$ be irrational with continued fraction expansion
$\alpha = [a_0; a_1, a_2, \ldots]$ where $a_0 = w = \floor{\alpha}$.
The convergents $p_k/q_k$ satisfy the recurrence
\begin{equation}\label{eq:cf-recurrence}
  p_{k+1} = a_{k+1} p_k + p_{k-1}, \qquad
  q_{k+1} = a_{k+1} q_k + q_{k-1},
\end{equation}
with $p_{-1} = 1$, $q_{-1} = 0$, $p_0 = w$, $q_0 = 1$.
The \emph{semi-convergents} between $p_{k-1}/q_{k-1}$ and
$p_{k+1}/q_{k+1}$ are
$\frac{p_{k-1} + j p_k}{q_{k-1} + j q_k}$ for
$j = 1, \ldots, a_{k+1} - 1$.

\subsection{Sturmian words}\label{ssec:sturmian}

The \emph{Sturmian word} of slope $1/\alpha$ is the binary sequence
$\mathbf{s} = (s_n)_{n \geq 0}$ with
$s_n = \floor{(n+1)/\alpha} - \floor{n/\alpha} \in \{0, 1\}$.
The letter $1$ indicates a \emph{rise} (height increases by one) and
$0$ a \emph{within-stair} (height stays).
Between convergent positions $p_{k-1}$ and $p_k$, the staircase
$\floor{x/\alpha}$ has exactly $q_k$ rises and
$p_k - q_k = (w-1)q_k + q_{k-1}$ within-stairs.

For a block with first partial quotient $a_1$, the Sturmian word
has pattern
\begin{equation}\label{eq:sturmian-pattern}
  \underbrace{0 \cdots 0}_{w}\; 1 \;
  \underbrace{0 \cdots 0}_{w-1}\; 1 \;
  \underbrace{0 \cdots 0}_{w-1}\; 1 \;
  \cdots \;
  \underbrace{0 \cdots 0}_{w-1}\; 1,
\end{equation}
with $q_1 = a_1$ rises.
The total length is $p_1 = wa_1 + 1$ columns.

\subsection{Transfer matrices and state vectors}

The \emph{state vector} $\mathbf{v}(c) = (v_0(c), v_1(c), \ldots,
v_N(c))$ at column $c$ counts the number of monotonic lattice paths
from $(1, s)$ to $(c, j)$ staying weakly below the staircase
$\floor{x/\alpha}$.
The column-by-column recurrence is:
\begin{align}
  v_j(c{+}1) &= \sum_{k=0}^{j} v_k(c)
  &&\text{(within-stair column)}, \label{eq:within}\\
  v_j(c{+}1) &= \sum_{k=0}^{j-1} v_k(c)
  &&\text{(rise column, row $j$ newly accessible)}.
  \label{eq:rise}
\end{align}
Starting from $v_j(0) = \delta_{j,s}$ for source height~$s$, the
\emph{block transfer matrix} $M \in \Z^{(N+1) \times (N+1)}$ is
defined by
\begin{equation}
  v_j(p_1) = \sum_{s=0}^{N} M[j, s]\, v_s(0).
\end{equation}

\subsection{Ballot numbers}

The ballot number $\B(p,q) = \frac{1}{p}\binom{p+q-1}{q}$ counts
lattice paths from $(1,0)$ to $(p,q)$ staying weakly below
$y = qx/p$ for coprime $p, q$.
It satisfies $\B(p, 0) = 1$ for all $p \geq 1$ and
$\B(1, q) = \delta_{q,0}$.

The \emph{generalized ballot formula} for height $j$ at position $p$
under a uniform staircase of step width $w$ is
\begin{equation}\label{eq:vlin}
  v_j(p) = \frac{p - wj}{p}\binom{p + j - 1}{j}.
\end{equation}
Note that $v_0(p) = 1$ for all $p$ and
$v_q(p) = \B(p, q)$ when $p = wq + 1$.

%======================================================================
\section{Uniform formula}\label{sec:uniform}
%======================================================================

\begin{theorem}[Uniform formula]\label{thm:uniform}
Let $\alpha = [w; a_1, a_2, \ldots]$ be irrational and let $p$
be a semi-convergent numerator between convergents $p_{k-1}$ and
$p_{k+1}$.
Then
\begin{equation}\label{eq:uniform}
  v_j(p) = \frac{p - wj}{p}\binom{p+j-1}{j}
  \qquad\text{for } j = 0, 1, \ldots, q_1.
\end{equation}
\end{theorem}

This is the rational ballot formula of Irving and
Rattan~\cite{IrvingRattan2009}, applied to the staircase
$\floor{qx/p}$ which agrees with $\floor{x/\alpha}$ on the first
$q_1 + 1$ heights.

\begin{proposition}[Two-term exactness]\label{prop:twoterm}
If $p/q = [w; a_1]$ is a rational with a two-term continued
fraction, then~\eqref{eq:uniform} is exact for \emph{all}
$j = 0, 1, \ldots, q$.
In particular, no correction is needed.
\end{proposition}

\begin{proof}
A two-term CF means the staircase $\floor{qx/p}$ has a single
Sturmian block with uniform step spacing.
The classical ballot formula applies at every height.
\end{proof}

\begin{corollary}\label{cor:correction-start}
Corrections to the uniform formula first appear at height
$j = q_1 + 1$ and are present precisely when $p/q$ has three or
more continued fraction terms (i.e., the staircase has non-uniform
step spacing).
\end{corollary}

%======================================================================
\section{Block transfer decomposition}\label{sec:transfer}
%======================================================================

\subsection{Toeplitz structure}

The \emph{Toeplitz matrix} for a block of width $p_1$ is
\begin{equation}\label{eq:toeplitz}
  T[j, s] = \binom{p_1 - 1 + j - s}{j - s},
\end{equation}
which counts all monotonic paths from $(0, s)$ to $(p_1, j)$
\emph{without} the staircase constraint.
This is the transfer matrix that would arise if all $p_1$ columns
were within-stairs (no rises).

\begin{proposition}\label{prop:toeplitz-exact}
For $j \leq d_0 - 1$, where $d_0 = q_0 + q_1 - 1$ is the input
state dimension (input rows $0, \ldots, d_0 - 1$), the transfer
matrix equals the Toeplitz matrix:
$M[j, s] = T[j, s]$.
\end{proposition}

\begin{proof}
A monotone path ending at height $j$ has height at most $j$ at
every column.
The staircase constraint is weakest at the start of the block,
where the accessible rows are $0, \ldots, d_0 - 1$.
Hence for $j \leq d_0 - 1$ every unconstrained path counted by
$T[j, s]$ already satisfies the constraint, and the two counts
agree.
\end{proof}

\subsection{The correction formula}

\begin{theorem}[Block transfer decomposition]\label{thm:transfer}
The block transfer matrix for a Sturmian block with parameters
$(d_0, w, p_1, q_1)$ decomposes as
\begin{equation}\label{eq:decomp}
  M[j, s] = T[j, s] - \Delta[j, s],
\end{equation}
where $\Delta[j, s] = 0$ for $j < d_0$ and, for
$j = d_0 + d$ with $0 \leq d \leq q_1 - 1$:
\begin{equation}\label{eq:R7}
  \Delta[d_0{+}d,\, s] = \sum_{m=1}^{d+1}
    v_{d-m+1}(p_1 - wm)\;\binom{d_0 - 2 + m(w{+}1) - s}{mw - 1}.
\end{equation}
Here $v_j(p)$ is the uniform ballot formula~\eqref{eq:vlin}.
\end{theorem}

\begin{remark}[Offset]\label{rmk:offset}
The offset $A = d_0 - 2$ in the binomial factor is an index
identity, not an independent phenomenon: for the row
$j_m = d_0 + m - 1$ newly accessible at the $m$-th rise,
\[
  \binom{mw - 1 + j_m - s}{j_m - s}
  = \binom{d_0 - 2 + m(w{+}1) - s}{mw - 1},
\]
so the binomial factor in~\eqref{eq:R7} is precisely the
\emph{free} path count into that row at the rise column
(Step~1 of the proof below).
In particular $A = d_0 - 2$ is universal --- independent of $w$,
$p_1$, and $q_1$.
\end{remark}

\begin{proof}[Proof of Theorem~\ref{thm:transfer}]
The proof tracks the correction vector $\boldsymbol{\delta}(c)$,
defined as $\delta_j(c) = T[j, s](c) - M[j, s](c)$, column by
column through the $q_1$ phases of the Sturmian block, and then
solves the resulting recursion in closed form by linearity and
self-similarity.

\emph{Phase structure.}
The block~\eqref{eq:sturmian-pattern} consists of $q_1$ phases,
each comprising one rise followed by within-stair columns:
\begin{itemize}
\item Phase $m$ ($1 \leq m \leq q_1 - 1$): rise at column
  $mw + 1$, then $w - 1$ within-stairs.
\item Phase $q_1$: rise at column $q_1 w + 1 = p_1$, no
  trailing within-stairs.
\end{itemize}
Before phase~$1$, there are $w$ initial within-stairs (columns
$1, \ldots, w$) which do not contribute to $\boldsymbol{\delta}$
since the correction starts at zero.

\emph{Step 1: correction dynamics.}
At a within-stair column, both the free evolution $T(c)$ and the
constrained evolution $M(c)$ obey the prefix-sum
recurrence~\eqref{eq:within}, so
$\boldsymbol{\delta}(c{+}1) = L\boldsymbol{\delta}(c)$, where $L$
denotes the prefix-sum operator.
At the rise of phase $m$ (the transition from column $mw$ to
$mw + 1$), the row $j_m = d_0 + m - 1$ (correction depth
$d = m - 1$) becomes accessible.
Existing rows again satisfy $\delta' = L\delta$, while the new
row satisfies, by~\eqref{eq:within} for $T$ and~\eqref{eq:rise}
for $M$,
\[
  \delta_{j_m}' = \sum_{k \leq j_m} T_k(mw)
   - \sum_{k < j_m} M_k(mw)
  = \sum_{k < j_m} \delta_k(mw) + t_m,
\]
with the scalar inhomogeneous term
\[
  t_m := T_{j_m}(mw) = \binom{mw - 1 + j_m - s}{j_m - s}
  = \binom{d_0 - 2 + m(w{+}1) - s}{mw - 1},
\]
the free path count into the newly accessible row
(Remark~\ref{rmk:offset}).
In operator form, each non-final phase acts as
\begin{equation}\label{eq:phase-op}
  \boldsymbol{\delta}^{(m+1)} = L^{w-1}\bigl(
    L(\boldsymbol{\delta}^{(m)} \oplus 0)
    + t_m\, \mathbf{e}_{m-1}\bigr),
\end{equation}
where $L^n[\delta](d) = \sum_{d'=0}^{d}
\binom{n-1+d-d'}{d-d'}\,\delta(d')$ is the $n$-fold prefix sum,
$\oplus 0$ denotes padding with a zero entry, and
$\mathbf{e}_{m-1}$ is the unit vector at position $m - 1$.
The final phase uses $L^0 = \mathrm{Id}$ instead of $L^{w-1}$.

\emph{Step 2: linearity.}
The dynamics~\eqref{eq:phase-op} are linear with a single scalar
injection $t_m$ per phase.
Starting from the empty vector, the final correction therefore
splits as
\[
  \delta_d^{(q_1)} = \sum_{m=1}^{q_1} t_m\, K_m(d),
\]
where $K_m(d)$ is the result of propagating a \emph{unit}
injection at depth $m - 1$ in phase $m$ through the remaining
phases.
The kernel $K_m$ is independent of $d_0$ and $s$, which enter
only through $t_m$; and since all operators are lower triangular,
$K_m(d) = 0$ for $d < m - 1$, so only the terms $m \leq d + 1$
contribute at depth $d$.

\emph{Step 3: self-similar kernel.}
After its injection, the unit at depth $m - 1$ undergoes the
following schedule of operations: $w - 1$ prefix sums (the rest of
phase $m$); then, for each of the phases $m + 1, \ldots, q_1 - 1$,
one append-and-prefix-sum followed by $w - 1$ prefix sums; finally
one append-and-prefix-sum (phase $q_1$).
These are verbatim the column recurrences of a staircase state
vector --- \eqref{eq:within} for a within-stair,
\eqref{eq:rise} for a rise --- with correction depth playing the
role of height.
Note that $K_m$ is \emph{not} the free composition
$L^{(q_1-m)w}$: the per-phase appending truncates the vector
(depth $d$ exists only from phase $d + 1$ on), and it is exactly
this truncation that re-creates the staircase constraint one
level down.
Relabel the injection column as $x = 1$ and set the relative
height $j = d - m + 1$, $q_{\mathrm{res}} = q_1 - m$, and
$P = w q_{\mathrm{res}} + 1 = p_1 - wm$.
The schedule has rises at $x = rw + 1$ for
$r = 1, \ldots, q_{\mathrm{res}}$, ending at $x = P$, with
within-stairs elsewhere.
Since $\lceil r P / q_{\mathrm{res}} \rceil = rw + 1$ for
$r = 1, \ldots, q_{\mathrm{res}}$, this is exactly the rise
schedule of the two-term staircase
$\floor{q_{\mathrm{res}}\, x / P}$.
Hence $K_m(d)$ equals the number of monotone lattice paths from
$(1, 0)$ to $(P, j)$ staying weakly below
$\floor{q_{\mathrm{res}}\, x / P}$, and by two-term exactness
(Proposition~\ref{prop:twoterm}),
\[
  K_m(d) = v_j(P) = v_{d-m+1}(p_1 - wm).
\]
Substituting into Step~2 yields~\eqref{eq:R7}.
\end{proof}

\begin{remark}[Independent verification]
Formula~\eqref{eq:R7} has additionally been verified symbolically
for $q_1 \leq 9$ with free parameters $w$, $d_0$, $s$, and against
explicitly computed block transfer matrices for
$\sqrt{2}$, $\sqrt{3}$, $\sqrt{5}$, $\sqrt{7}$, $\sqrt{37}$,
$\varphi$, $e$, and $1 + \pi/10$.
\end{remark}

\begin{example}\label{ex:sqrt5}
For $\alpha = \sqrt{5}$, $w = 2$, $a_1 = 4$, $p_1 = 9$, $q_1 = 4$,
$d_0 = 4$, the correction matrix $\Delta$ has non-zero entries at
rows $j = 4, \ldots, 7$ (i.e., $d = 0, \ldots, 3$); the first two
are
\begin{align*}
  \Delta[4, s] &= \binom{5 - s}{1},\\
  \Delta[5, s] &= 5\binom{5 - s}{1} + \binom{8 - s}{3}.
\end{align*}
These match~\eqref{eq:R7} with $d = 0$: one term ($m = 1$,
$v_0(7) = 1$, $\binom{5 - s}{1}$); and $d = 1$: two terms
($m = 1$: $v_1(7) = 5$, $\binom{5 - s}{1}$; $m = 2$:
$v_0(5) = 1$, $\binom{8 - s}{3}$).
For instance, at $s = 0$: $\Delta[4, 0] = 5$ and
$\Delta[5, 0] = 25 + 56 = 81$, matching the directly computed
block transfer matrix.
\end{example}

%======================================================================
\section{Hypergeometric classification}\label{sec:hypergeom}
%======================================================================

\subsection{The $w = 1$ collapse: Rothe--Hagen identity}

\begin{theorem}[Rothe--Hagen collapse]\label{thm:rothe-hagen}
For $w = 1$, the multi-term sum~\eqref{eq:R7} collapses to a
single binomial:
\begin{equation}\label{eq:w1-collapse}
  \Delta_{w=1}[d_0 + d,\, s] = \binom{d_0 + p_1 + d - 1 - s}{d}.
\end{equation}
\end{theorem}

\begin{proof}
For $w = 1$, $p_1 = q_1 + 1$, and the sum~\eqref{eq:R7} becomes
\[
  \sum_{m=1}^{d+1} v_{d-m+1}(p_1 - m)\,\binom{d_0 - 2 + 2m - s}{m-1}.
\]
Substituting $v_j(p) = \frac{p-j}{p}\binom{p+j-1}{j}$ and
simplifying, this is an instance of the Rothe--Hagen identity
\cite{Rothe1793, Hagen1891}
\[
  \sum_{k=0}^{n} \frac{a}{a+bk}\binom{a+bk}{k}
  \frac{c}{c+(n-k)b}\binom{c+(n-k)b}{n-k}
  = \frac{a+c}{a+c+bn}\binom{a+c+bn}{n}
\]
with parameters $b = 2$, evaluated at the appropriate $a$, $c$,
$n$.
The collapsed form~\eqref{eq:w1-collapse} has been verified
symbolically for $d = 0, \ldots, 5$ with free $d_0$, $p_1$, $s$.
\end{proof}

\subsection{General $w$: Saalsch\"utzian hypergeometric}

\begin{theorem}[Hypergeometric classification]%
\label{thm:hypergeom}
For general $w \geq 1$, the correction
sum~\eqref{eq:R7} at $s = 0$ equals, up to a prefactor, a
terminating $1$-balanced (Saalsch\"utzian) generalized
hypergeometric function:
\[
  \Delta[d_0{+}d,\, 0] = \mathrm{prefactor}(w, d_0, p_1, d)
  \cdot \pFq{2w+2}{2w+1}\!\left(
    \begin{matrix}
      \text{upper parameters} \\
      \text{lower parameters}
    \end{matrix}
  ;\, 1\right).
\]
The classification by $w$:
\begin{center}
\begin{tabular}{clll}
\toprule
$w$ & Type & Structure & Status \\
\midrule
$1$ & $\pFq{4}{3}(1)$ & Collapses (Pfaff--Saalsch\"utz) & $= \binom{\cdot}{\cdot}$ \\
$2$ & $\pFq{6}{5}(1)$ & Irreducible Saalsch\"utzian & Closed form \\
$3$ & $\pFq{8}{7}(1)$ & Irreducible Saalsch\"utzian & Closed form \\
$w$ & $\pFq{2w+2}{2w+1}(1)$ & Irreducible for $w \geq 2$ & Closed form \\
\bottomrule
\end{tabular}
\end{center}
All instances satisfy the $1$-balanced condition
$\sum(\text{upper}) - \sum(\text{lower}) + 1 = 0$.
\end{theorem}

\begin{proof}
Mathematica's \texttt{Sum} evaluates the
sum~\eqref{eq:R7} at $s = 0$ in terms of
\texttt{HypergeometricPFQ} for each fixed $w$.
The parameter lists are explicit functions of $d_0$, $p_1$, $d$.
The $1$-balanced property is verified symbolically.
For $w = 1$, the $\pFq{3}{2}(1)$ (the leading term is
$\pFq{4}{3}$ in the full parametrization, but after absorbing
the prefactor it reduces to $\pFq{3}{2}$) satisfies the
Pfaff--Saalsch\"utz identity, which forces collapse to a single
binomial coefficient.
For $w \geq 2$, \texttt{FunctionExpand} does not reduce the
$\pFq{2w+2}{2w+1}$, and no Whipple-type transformation applies.
\end{proof}

\begin{remark}[Pole regularization]
At integer specializations of $d_0$ and $p_1$, some lower
parameters of the $\pFq{2w+2}{2w+1}$ become non-positive integers,
creating Pochhammer poles.
The original sum~\eqref{eq:R7} remains well-defined; the
hypergeometric representation requires a limit.
All $27$ pole cases tested match exactly after regularization.
\end{remark}

%======================================================================
\section{Self-similar hierarchy}\label{sec:selfsim}
%======================================================================

\subsection{Universality across CF levels}

\begin{theorem}[Self-similar correction]\label{thm:selfsim}
The correction formula~\eqref{eq:R7} applies at every continued
fraction level $k$ with parameters $(d_{0,k}, w, p_k, q_k)$,
where $d_{0,k}$ is the state dimension at level $k$.
\end{theorem}

\begin{proof}
The derivation of~\eqref{eq:R7} depends only on the Sturmian block
structure: $q_k$ rises of type~\eqref{eq:rise} interspersed with
within-stairs of type~\eqref{eq:within}, in the
pattern~\eqref{eq:sturmian-pattern} with local parameters.
Since this structure is universal to all CF levels --- each
level~$k$ block has a Sturmian pattern with parameters
$(p_k, q_k, w)$ --- the same derivation applies.
The offset $A_k = d_{0,k} - 2$ is universal by
Remark~\ref{rmk:offset}.
\end{proof}

\subsection{Inclusion--exclusion composition}

At continued fraction level $k + 1$, the block transfer is a
product of $a_{k+1}$ sub-blocks:
\begin{equation}\label{eq:product}
  M_{k+1} = \prod_{i=1}^{a_{k+1}} (T_i - \Delta_i),
\end{equation}
where each factor $T_i - \Delta_i$ is the transfer matrix of a
level-$k$ sub-block with correction given
by~\eqref{eq:R7}.

\begin{theorem}[Terminating expansion]\label{thm:ie}
The inclusion--exclusion expansion of~\eqref{eq:product},
\begin{equation}\label{eq:ie-expansion}
  M_{k+1} = \sum_{S \subseteq \{1, \ldots, a_{k+1}\}}
  (-1)^{|S|} \prod_{i \in S} \Delta_i \prod_{i \notin S} T_i,
\end{equation}
terminates at order $|S| = a_{k+1}$.
Moreover, order $n$ first contributes at correction depth
$d = (n-1)q_k$, so for small $d$ only low-order terms are
needed.
\end{theorem}

\begin{proof}
The sum~\eqref{eq:ie-expansion} is a finite expansion of a product
of $a_{k+1}$ matrices and terminates by construction at
$|S| = a_{k+1}$.
The activation depth follows from the support of $\Delta_i$:
each correction $\Delta_i$ is non-zero only at rows
$j \geq d_{0,k} + 2$, contributing a dimension shift of $q_k$
per correction factor.
\end{proof}

\begin{example}
For $\alpha = \sqrt{5}$ ($a_2 = 4$), the level-$2$ block decomposes
into $4$ sub-blocks.
Order~$0$ (pure Toeplitz) gives the leading approximation;
order~$1$ corrects starting at $d = 0$;
order~$2$ at $d = 4$; order~$3$ at $d = 8$; order~$4$ at $d = 12$.
Each successive order improves accuracy by approximately two orders
of magnitude.
For $\alpha = \pi$ ($a_2 = 15$), the expansion has $15$ terms.
\end{example}

%======================================================================
\section{Rotation decomposition}\label{sec:rotation}
%======================================================================

The Sturmian word~\eqref{eq:sturmian-pattern} assumes a
\emph{canonical} starting position (first column is a within-stair).
At non-canonical positions within the staircase, the block begins
with $r \geq 1$ rises before the first within-stair.

\begin{definition}
The \emph{rotation parameter} at position $x$ in the staircase
$\floor{x/\alpha}$ is
$r(x) = \#\{\text{consecutive rises starting at column } x\}$.
\end{definition}

\begin{theorem}[Rotation decomposition]\label{thm:rotation}
For $w = 1$ and rotation $r$, the block transfer matrix decomposes
as
\begin{equation}\label{eq:rotation}
  M_r = \operatorname{sturmianBlock}(d_0 + r,\, p - r,\, q - r)
  \cdot \operatorname{riseProduct}(d_0,\, r),
\end{equation}
where $\operatorname{sturmianBlock}$ applies~\eqref{eq:R7} with
shifted parameters and $\operatorname{riseProduct}$ is the
composition of $r$ consecutive rise matrices from
dimension~$d_0$.
\end{theorem}

This has been verified for $\sqrt{2}$, $\sqrt{3}$, and $\varphi$
at continued fraction levels $1$ through~$6$ with exact match.

\begin{remark}
The rotation $r$ is the only non-symbolic component in the
closed-form transfer: it requires $O(\log p)$ evaluations of
$\floor{x/\alpha}$.
For quadratic irrationals with periodic CF, $r$ follows a
quasi-periodic pattern but has no known closed form.
\end{remark}

%======================================================================
\section{The $w$-reduction kernel}\label{sec:wreduction}
%======================================================================

For $w \geq 2$, the irreducible Saalsch\"utzian
$\pFq{2w+2}{2w+1}(1)$ of Section~\ref{sec:hypergeom} can be
related back to the $w = 1$ case via a convolution.

\begin{theorem}[$w$-reduction]\label{thm:wreduction}
For any $w \geq 2$, the correction at $s = 0$ satisfies
\begin{equation}\label{eq:convolution}
  \Delta_w[d, 0] = \sum_{j=0}^{d} K_w[j;\, d_0, q_1]\,
  \Delta_1[d - j, 0],
\end{equation}
where $\Delta_1[d, 0] = \binom{d_0 + q_1 + d}{d}$ is the
Rothe--Hagen form and $K_w$ is a polynomial kernel with:
\begin{enumerate}
\item $\deg K_w = (w-1)(d_0 + q_1) - w$;
\item $K_w[0] = \binom{d_0 + w - 1}{w - 1}$ (universal, independent of $q_1$);
\item Leading coefficient: $[q_1^d]\, K_w[d]
  = \binom{d_0 + w - 1}{w - 1}\cdot\frac{(w-1)^d}{d!}$.
\end{enumerate}
\end{theorem}

The kernel $K_w$ has been computed explicitly for $w = 2, 3, 4$
and verified across $\sqrt{5}$, $\sqrt{6}$, $\sqrt{7}$ for
$d_0 = 3, \ldots, 10$ with exact match in all $24$ test cases.

\begin{remark}
The leading coefficient formula implies rapid decay:
$(w-1)^d/d!$ decreases factorially, so the kernel concentrates
near $j = 0$.
This explains why the $w \geq 2$ corrections are well-approximated
by the $w = 1$ (Rothe--Hagen) correction with a multiplicative
factor of $\binom{d_0 + w - 1}{w - 1}$.
\end{remark}

%======================================================================
\section{Level scaling for quadratic irrationals}%
\label{sec:scaling}
%======================================================================

For $w = 1$ and quadratic irrationals $\alpha = \sqrt{D}$, the
continued fraction is eventually periodic, and the block corrections
at successive levels exhibit a clean scaling law.

\begin{theorem}[Fundamental unit scaling]\label{thm:scaling}
For $w = 1$, $\alpha = \sqrt{D}$ with periodic CF, the full-block
correction at level $k$ satisfies
\begin{equation}\label{eq:scaling}
  C_k(d, s) = A_d\, \varepsilon^k + B_d\, \bar{\varepsilon}^k + H_d,
\end{equation}
where $\varepsilon$ is the fundamental unit of
$\mathbb{Q}(\sqrt{D})$, $\bar{\varepsilon}$ is its algebraic
conjugate, and $A_d$, $B_d$, $H_d$ depend on $d$, $s$, and the
CF parameters.
\end{theorem}

\begin{remark}
This is \emph{not} a new method for computing $\varepsilon$ (the CF
is faster).
The significance is combinatorial: it provides a new interpretation
of the fundamental unit as the scaling eigenvalue of lattice path
corrections under continued fraction refinement.
\end{remark}

The recursion is triangular: the Toeplitz convolution at row
$d_0 + D$ depends only on corrections at rows $0, \ldots, D$, so
the recursion closes for any fixed depth.
At $d = 0$, the corrections alternate $1, 0, 1, 0, \ldots$
(eigenvalues $\pm 1$).

%======================================================================
\section{Open problems}\label{sec:open}
%======================================================================

\begin{enumerate}
\item \textbf{Rotation decomposition for $w \geq 2$.}
  Theorem~\ref{thm:rotation} is proved for $w = 1$.
  Extending it to general $w$ requires a rotation-aware version of
  the multi-term sum~\eqref{eq:R7}.

\item \textbf{Sub-leading kernel coefficients.}
  The three universal formulas for $K_w$ (degree, $K_w[0]$, leading
  coefficient) are proved, but the sub-leading $q_1$-coefficients
  have no known universal closed form.

\item \textbf{SL$(2, \mathbb{Z})$ connection.}
  The sub-block factorization of Section~\ref{sec:selfsim} mirrors
  the continued fraction matrix recursion
  $\bigl(\begin{smallmatrix} p_{k+1} & p_k \\ q_{k+1} & q_k
  \end{smallmatrix}\bigr) =
  \bigl(\begin{smallmatrix} a_{k+1} & 1 \\ 1 & 0
  \end{smallmatrix}\bigr)
  \bigl(\begin{smallmatrix} p_k & p_{k-1} \\ q_k & q_{k-1}
  \end{smallmatrix}\bigr)$.
  A precise correspondence with the SL$(2, \Z)$ decomposition of
  the transfer matrix is unknown.

\item \textbf{Irrationality measure connection.}
  The combinatorial deficit
  $\delta_j(p) = v_j(p) - v_j^{\mathrm{lin}}(p)$ encodes continued
  fraction structure.
  Whether this connects to irrationality measures or Diophantine
  approximation quality is open.
\end{enumerate}

%======================================================================
\section*{Acknowledgments}
%======================================================================

Computational experiments were performed in Wolfram Language using
the Orbit paclet
(\url{https://github.com/popojan/orbit}).

\bibliographystyle{plain}
\bibliography{references}

\end{document}
